\documentclass[10pt,a4paper]{amsart}
\usepackage[margin=19mm]{geometry}
\usepackage{amsmath,amssymb,mathtools}
\usepackage[scr=rsfs,cal=euler]{mathalfa}
\usepackage{xfrac}
\usepackage[unicode,hidelinks,bookmarks=false]{hyperref}
\newcommand{\ie}{i.e.\ }
\newcommand{\cf}{cf.\ }
\newcommand{\resp}{respectively\ }
\newcommand{\Subsection}{\subsection}
\providecommand{\nobreakdash}{\nobreak\hbox{-}}
\setlength{\emergencystretch}{2em}
\multlinegap0pt
\newtheorem{thm}{Theorem}
\newtheorem{lemma}[thm]{Lemma}
\newtheorem{prop}[thm]{Proposition}
\newtheorem{cor}[thm]{Corollary}
\newtheorem{exam}[thm]{Example}
\newtheorem{rem}[thm]{Remark}
\newtheorem{defi}[thm]{Definition}
\def\g{\gamma}
\def\ds{\displaystyle}
\setlength{\topmargin}{-0.55cm} \setlength{\oddsidemargin}{.40cm}
\setlength{\evensidemargin}{.40cm} \setlength{\textheight}{23.5cm}
\setlength{\textwidth}{15cm}
\newcommand{\sign}{\text{\rm sign}}
\begin{document}
\title[Geometric optimization problems generated by plane curves]{Geometric optimization problems generated by plane curves}
\author{Petar Kenderov P. Kenderov; Acad. G. Bonchev 8, 1113 Sofia, Bulgaria špace 1mm; Petar Kenderov, Oleg Mushkarov, Nikolai Nikolov}
\date{}
\maketitle
\noindent\emph{Source and licence.} Geometric optimization problems generated by plane curves. Version arXiv:2608.10128v1 (CC BY 4.0). This material is distributed under the Creative Commons Attribution 4.0 International licence, http://creativecommons.org/licenses/by/4.0/, as recorded in the arXiv OAI licence metadata for this exact version and shown on the abstract page https://arxiv.org/abs/2608.10128v1. Full rights notice: licenses/arxiv-2608.10128-CC-BY-4.0.txt.\par\smallskip
\noindent\textit{Editorial scope.} Counted unit: the original \texttt{thm} environment \texttt{A} at source lines 444--459 together with its complete proof at lines 461--600; the delivered document keeps source lines 216--601 so that every referenced label and every citation resolves inside it. Native editable source is the paper's own concatenated TeX; the AMS wrapper is editorial formatting only. No publisher class, upstream script or unrelated artwork is redistributed.
\begin{equation}\label{normala}
(x - p(t_0))p\,'(t_0) + (y - q(t_0))q\,'(t_0) = 0.
\end{equation}

Recall that the signed curvature $\kappa_{\g}(t)$ at some point $\g(t)=(p(t),q(t))$, $t \in \Delta$, of a $C^2$-smooth
plane curve  $\g$  is given by

\begin{equation}\label{1}
\kappa_{\g}(t)=\frac{p\,'(t)q\,''(t)-p\,''(t)q\,'(t)}{(p\,'(t)^2+q\,'(t)^2)^{3/2}} .
\end{equation}

%When it is necessary to stress the role of the
%point $\g(t)$ we use  also the symbol $\kappa_{\g}(\g(t))$ instead of
%$\kappa_{\g}(t)$.

If $s$ is the arc-length parameter of $\g$, then
\begin{equation}\label{2}
\kappa_{\g}(s)=p\,'(s)q\,''(s)-q\,'(s)p\,''(s)
\end{equation}
since $p\,'(s)^2+q\,'(s)^2=1$. If
$\kappa_{\g}(s_0)\neq 0$ then the center of
curvature $C_{\g(s_0)}$ of $\g$ at $\g(s_0)$ is the point
\begin{equation}\label{3}
C_{\g(s_0)}=(p(s_0)-\frac{q\,'(s_0)}{\kappa_{\g}(s_0)} \ , \
q(s_0)+\frac{p\,'(s_0)}{\kappa_{\g}(s_0)}).
\end{equation}
%In vector form, $C_{\g(s_0)} - \g(s_0) = \frac{1}{\kappa_{\g}(s_0)}N(s_0)$ .
The center of curvature $C_{\g(s_0)}$ lies on the normal $n_{\g(s_0)}$ to $\g$ at
$\g(s_0)$ and the distance between $\g(s_0)$
and $C_{\g(s_0)}$ is equal to $|\frac{1}
{\kappa_{\g }(s_0)}|$ because $N(s_0)$ is a vector of lenght one.


Let $\g_1(t)$, $t \in \Delta_1$, and $\g(s)$, $s \in \Delta$, be $C^1$-smooth curves
with  $s$ being the arc-length parameter of  $\g$.  Consider the function
\begin{equation}\label{G}
G(t,s)=(p_1(t)-p(s))p\,'(s)+(q_1(t)-q(s))q\,'(s)
\end{equation}
defined in $\Delta_1 \times \Delta$.

The next simple statement summarizes some of the relations between the function $G$ and its partial derivatives
 $G'_t$ and $G'_s$ on the one hand, and the geometric essence of the problems considered in this paper on the other.

\begin{lemma}\label{l1}
(i) A pair $(t_0, s_0)$ is a solution to the equation $G(t,s)=0$ if and only if $\g_1(t_0) \in n_{\g(s_0)}$.

(ii) A pair $(t_0, s_0)$ is a solution to the system of two equations $G(t,s)=0$,
$G'_t(t,s)=0$ if and only if
$n_{\g(s_0)} = t_{\g_1(t_0)}$.

(iii)  If $\g(s)$, $s \in \Delta$, is a $C^2$-smooth curve then a pair $(t_0, s_0)$ is a solution to the system of two equations $G(t,s)=0$,
$G'_s(t,s)=0$ if and only if $\kappa_{\g}(s_0) \neq 0$ and $\g_1(t_0) = C_{\g(s_0)}$.

\end{lemma}

\begin{proof}
Item (i) follows trivially from
%the regularity of $\g$ and
the definition of the normal $n_{\g(s_0)}$ in terms of equation (\ref{normala}).

The relation
$ G'_t(t_0,s_0)= p\,'_1(t_0)p\,'(s_0)+q\,'_1(t_0)q\,'(s_0)=0 $ shows that the non-zero tangent vectors
$\g\,'_1(t_0)=(p\,'_1 (t_0), q\,'_1(t_0))$ and
$\g\,'(s_0)=(p\,'(s_0), q\,'(s_0))$ are orthogonal to each other. Combined with the fact that $\g_1(t_0) \in n_{\g(s_0)}$, this yields (ii).

To prove (iii), we note that
\begin{equation}\label{4}
G'_s(t_0,s_0)=
(p_1(t_0)-p(s_0))p\,''(s_0)+(q_1(t_0)-q(s_0))q\,''(s_0)-1
\end{equation}
 because $p\,'(s_0)^2+ q\,'(s_0)^2=1$.
 If $G'_s(t_0,s_0)=0$, equation (\ref{4}) implies that the vector
 $\g\,''(s_0)=(p\,''(s_0), q\,''(s_0)) \neq (0,0)$.
 Hence,
 $\kappa_{\g}(s_0) \neq 0$ and the point $C_{\g(s_0)}$ exists.

Taking into account the Frenet-Serret formula
\begin{equation}\label{5}
(p\,''(s),q\,''(s)) = \kappa_{\g}(s)(-q\,'(s), p\,'(s))
\end{equation}
we get
\begin{equation}\label{6}
G'_s(t_0,s_0)=
\kappa_{\g}(s_0)[(p_1(t_0)-p(s_0))(-q\,'(s_0))+(q_1(t_0)-q(s_0))p\,'(s_0)]-1
\end{equation}
Having in mind (\ref{3}), i.e. that $C_{\g(s_0)} - \g(s_0) = \frac{1}{\kappa_{\g}(s_0)}N(s_0)$,
this equation can be rearranged in terms of the scalar product of vectors as

\begin{equation}\label{7}
G'_s(t_0,s_0)=\kappa_{\g}(s_0) \langle(\g_1(t_0) - \g(s_0)), N(s_0)\rangle - 1= \kappa_{\g}(s_0) \langle(\g_1(t_0) - C_{\g(s_0)}), N(s_0)\rangle.
\end{equation}
This suffices to complete the proof of item (iii) because the vectors $N(s_0)$, $\overrightarrow{C_{\g(s_0)}\g_1(t_0)}$, and
$\overrightarrow{\g(s_0)\g_1(t_0)}$ are colinear if and only if $G(t_0,s_0)=0$ (follows from item (i) ).
\end{proof}

Next, we provide sufficient conditions under which a curve
$\g_1(t), t \in \Delta_1$, can be locally reparametrized by the  arc-length parameter $s$ of another curve $\g(s), s \in \Delta$, and,
vice-versa, the curve $\g(s)$ can be locally reparametrized by the parameter $t$ of the curve $\g_1(t)$.

\begin{prop} \label{nnorm1}
 Let $\g_1(t)$, $t \in \Delta_1$, and $\g(s)$,
 $s \in \Delta$, be  $C^1$-smooth curves where  $s$ is the arc-length parameter of
 $\g$. Let $A_1= \g_1(t_0)$ and $A= \g(s_0)$ be such that $A_1 \in n_A$.

(i) If $n_A$ is not tangent to $\g_1$ at $A_1$,
then there exist an open interval $U \subset \Delta$ containing $s_0$, and a unique $C^1$-smooth function $t_1(s)$ defined in $ U$ with values in
$\Delta_1$ such that $t_1(s_0) = t_0$ and for every
$s \in U$ the point $\g_1(t_1(s))$ belongs to the normal $n_{\g(s)}$.

(ii) If $\g(s)$, $s \in \Delta$, is a $C^2$-smooth curve and either $\kappa_{\g}(A)=0$ or $A_1\neq C_A$, then there exist an open interval $V \subset \Delta_1$ containing $t_0$, and a unique $C^1$-smooth function $s_1(t)$ defined in $V$ with values in
$\Delta$ such that $s_1(t_0) = s_0$ and for every
$t \in V$ the point $\g_1(t)$ belongs to the normal $n_{\g(s_1(t))}$.
\end{prop}

\begin{proof}
(i) As $A_1 \in n_A$ we get from item (i) of Lemma \ref{l1} that
$G(t_0,s_0)=0$. Since $n_A$ is not tangent to $\g_1$ at $A_1$ we get from item (ii) of Lemma \ref{l1} that $G'_t(t_0,s_0) \neq 0$.
The implicit function theorem provides us with an open interval $U$, $s_0 \in U \subset \Delta$, and with a unique continuously differentiable function $t_1(s)$ defined in $U$ with values in
$\Delta_1$ such that $t_1(s_0)=t_0$ and $G(t_1(s), s) =0$ for every $s \in U$. This means,  $\g_1(t_1(s))$ belongs to the normal $n_{\g(s)}$ for every $s \in U$. The proof of item (i) is completed.

The proof of (ii) is similar. If $\kappa_{\g}(A)=0$ or $A_1\neq C_A$, we get from the proof of item (iii), Lemma \ref{l1},  that in both cases $G'_s(t_0,s_0) \neq 0$. Using the implicit function theorem again, we get an open interval $V $, $t_0 \in V \subset \Delta_1$, and a unique continuously differentiable function $s_1(t)$ defined in $V$ with values in $\Delta$ such that $s_1(t_0)=s_0$ and
$G(t,s_1(t)) =0$ for every $t\in V$.
This means, $\g_1(t)$ belongs to the normal $n_{\g(s_1(t))}$ for every $t \in V$. This completes the proof of (ii).

\end{proof}

\begin{prop}\label{nnorm2}
 Let $\g_1(t)$, $t \in \Delta_1$, and $\g(s)$,
 $s \in \Delta$, be  $C^1$-smooth curves (with  $s$ being the arc-length parameter of
 $\g$). Let the points $A_1= \g_1(t_0)$ and $A= \g(s_0)$ be such that $A_1 \in n_A$.

If $n_A$ is not tangent to $\g_1$  at $A_1$, then  the following properties are equivalent:

(a) The derivative $t_1'(s_0) \neq 0$ where
$t_1: U \rightarrow \Delta_1$ is the $C^1$-smooth function from item (i) of Proposition
\ref{nnorm1};

(b) There exists an open subinterval $U'$, $s_0 \in U'\subset U$, on which $t_1$ is a difeomorphism (i.e. the restriction of $t_1$ on $U'$ is invertible and its inverse $t_1^{-1}: t(U') \rightarrow U' $  is a  $C^1$-smooth mapping).

If $\g$ is a $C^2$-smooth curve, then each of the items (a) and (b) is equivalent to the next property

(c)  Either $C_A$ does not exist (i.e.
$\kappa_{\g}(A)=0$) or $A_1\neq C_A$.

\end{prop}
\begin{proof} If $t_1'(s_0) \neq 0$  then $t_1'(s) \neq 0$ for all $s$ from some sub-interval $U'$, $s_0 \in U'\subset U$. This means, the mapping $t_1:U' \to t_1(U')$ is strictly monotone in $U'$ and, therefore, invertible. Moreover, since $t_1$ is continuously differentiable with non-vanishing derivative, its inverse
$t_1^{-1}:t_1(U') \rightarrow U'$ is continuously differentiable as well, and has a non-vanishing derivative, i.e., (a) and (b) are equivalent.

We show now that (a) and (c) are equivalent if $\g$ is a $C^2$-smooth curve.
%Suppose  $t_1'(s_0) \neq 0$.
By assumption,
$n_A$ is not tangent to $\g_1$ at $A_1$. Hence, by item (ii) of Lemma \ref{l1} we have $G'_t(t_0,s_0) \neq 0$. Consider the function $t_1(s)$ from item (i) of Proposition \ref{nnorm1} and differentiate the identity  $G(t_1(s), s) =0$. We get, for $s=s_0$, the equation

\begin{equation}\label{9}
G'_t(t_1(s_0),s_0)t_1'(s_0) + G'_s(t_1(s_0),s_0)=G'_t(t_0,s_0)t_1'(s_0) + G'_s(t_0,s_0) = 0.
\end{equation}

It follows that $t'_1(s_0) \neq 0$ if and only if $G_s'(t_0,s_0) \neq 0$, which, by virtue of item (iii) of Lemma \ref{l1}, is equivalent to  (c).
\end{proof}
The next statement is a ``mirror version'' of the previous one, and its proof is omitted.


\begin{prop} \label{nnorm3}
Let $\g_1$ be a $C^1$-smooth curve, $\g$ be a $C^2$-smooth curve, $A =\g(s_0)$, and
$A_1=\g_1(t_0)$.
If either $\kappa_{\g}(s_0)=0$ or $A_1\neq C_A$ then  the following properties are equivalent:

(a) The derivative $s_1'(t_0) \neq 0$ where  $s_1: V \rightarrow \Delta$ is the $C^1$-smooth function from item (ii) of Proposition
\ref{nnorm1};

(b) there ixists an open subinterval $V'$, $s_0 \in V'\subset V$, on which $s_1$ is a difeomorphism (i.e. the restriction of $s_1$ on $V'$ is invertible and its inverse $s_1^{-1}: s(V') \rightarrow V' $  is a  $C^1$-smooth mapping);

(c) $n_A$ is not  tangent to $\g_1$  at $A_1$

\end{prop}

\begin{cor} \label{homo}
 Let $\g_1(t)$, $t \in \Delta_1$ be a $C^1$-smooth curve, $\g(s)$,  $s \in \Delta$, be  $C^2$-smooth curve and let $A_1= \g_1(t_0)$ and $A= \g(s_0)$ be such that $A_1 \in n_A$.

If either $\kappa_{\g}(A) = 0$ or $A_1\neq C_A$, and $n_A$ is not tangent to $\g_1$ at $A_1$,  then there exists an interval $U''$,  $s_0 \in U'' \subset \Delta$, such that the restriction of $t_1$ on $U''$ and the restriction of $s$ on $t_1(U'')$ are inverse to each other.
\end{cor}

\begin{proof}
Consider the inverse mapping $t_1^{-1}:t(U') \rightarrow U'$ from Proposition \ref{nnorm2}. Its restriction on $V' \cap t(U')$, where $V'$ is from Proposition \ref{nnorm3}, is a second implicit function (the first being the restriction of $s_1$ on this set). Uniqueness of the implicit function implies that $s_1$ and $t_1^{-1}$ coincide on the set $V' \cap t(U')$.

\end{proof}

We will also need the following simple observation.

\begin{lemma}\label{lem} Let $\g_1(t)$, $t \in \Delta_1$, and $\g_2(\tau)$, $ \tau \in \Delta_2$, be $C^1$-smooth curves. Let $t_0 \in W$ where
$W$ is an open subset of $\Delta_1$ and let
$u:W \to \Delta_2$ be a $C^1$-smooth function such that $\g_1(t_0) \neq \g_2(u(t_0))$ and $u\,'(t_0)=0$.
Set $d(t)= |\g_1(t)\g_2(u(t))|$ for any $t \in W$.
Then $d'(t_0)=0$ if and only if the line
$\g_1(t_0)\g_2(u(t_0))$ is the normal to $\g_1$ at $\g_1(t_0)$.
\end{lemma}

\begin{proof} Since $u'(t_0) =0$, we get for the derivative  at $t_0$ of the function
$$d(t)^2 = \langle \g_1(t) - \g_2(u(t)), \g_1(t) - \g_2(u(t))\rangle$$
the expression
$$ d(t_0) d'(t_0) = \langle \g_1(t_0) - \g_2(u(t_0)), \g_1'(t_0) \rangle.$$
As $d(t_0) \neq 0$ we see that
 $d'(t_0)=0$ if and only if the non-zero vector
 $\g_1(t_0) - \g_2(u(t_0))$ is perpendicular to the non-zero tangent vector $\g\,'_1(t_0)$.
\end{proof}



\section{Main results}\label{mr}

Let $\g_1$ and $\g_2$ be $C^1$-smooth curves in the plane.
In this section, we obtain geometric information about the points of local extremum of the distance function
$d(A_1,A_2)=|A_1A_2| , A_1\in \g_1 , A_2\in \g_2$, $A_1 \neq A_2$, under the following constraints:

(i) The lines $A_1A_2$ are normal to a given $C^2$-smooth curve $\g$.

(ii) The lines $A_1A_2$ pass through a given point $O$.

(iii) The lines $A_1A_2$ are tangent to a given $C^1$-smooth curve $\g$.

(iv) The lines $A_1A_2$ intersect a given $C^1$-smooth curve $\g$.

To distinguish these four cases, we denote the corresponding
distance functions by $d_{n_{\g}}, d_O \ , d_{t_{\g}}, $ and
$d_{\g}$.

We first consider the function $d_{n_{\g}}$.


\begin{thm}\label{A} Let $\g_1(t)$,
$t\in \Delta_1$, and $ \g_2(\tau)$,
$\tau \in \Delta_2$,  be $C^1$-smooth curves
and $\g(s)$, $s \in \Delta$, be a $C^2$-smooth curve with sign curvature
$\kappa_{\g}$. Let $A_1\in \g_1$ and $A_2\in \g_2$ be such that
$A_1\neq A_2$, $A_1A_2=n_A$ for some $A \in \g$, and the function $d_{n_{\g}}$ has a local extremum at the triple $(A, A_1,A_2)$.


(i) If $\kappa_{\g}(A)\neq 0$, then the normals $n_{A_1} ,
n_{A_2}$ and the line $n_{C_A}$ through the center of curvature
$C_A$ of $\g$ and perpendicular to $A_1A_2$ either intersect at
one point, or are parallel.

(ii) If $\kappa_{\g}(A)= 0$, then the normals $n_{A_1}$ and
$n_{A_2}$ are parallel.
\end{thm}

\begin{proof}To prove (i), we consider the following three options for the line $A_1A_2$ that cover all possible cases:

\emph{1. $A_1A_2$ is tangent to $\g_1$ and  $\g_2$
at $A_1$ and  $A_2$ correspondingly.}

\emph{2. $A_1A_2$ is not tangent to $\g_1$ and
$\g_2$ at $A_1$ and $A_2$ correspondingly. }

\emph{3. $A_1A_2$ is tangent to $\g_1$ at $A_1$ and not tangent to $\g_2$ at $A_2$. }

In the case \emph{1} the normals $n_{A_1}$ and $n_{A_2}$ are parallel to
the line $n_{C_A}$ and there is nothing to prove.

In case \emph{2}, we will prove that the normals
$n_{A_1} , n_{A_2}$
and the line $n_{C_A}$ intersect at one point.
We split case \emph{2} into two separate sub-cases:

(a) $C_A$ is different from the points $A_1$ and $A_2 $ (this is the so-called ``generic case'');

(b) $C_A$ coincides with one of the two different points $A_1$ and $A_2$.

Consider first the sub-case (a).
Let $s$ be the arc-length parameter of $\g$ with $A=\g(0)$
and set $\g(s)=(p(s), q(s)) \ , \ T(s)=\g\,'(s)=(p\,'(s),q\,'(s))$ and
$N(s)= (-q\,'(s), p\,'(s)$. As already mentioned above,  $T(s)$ is the unit tangent vector
and $N(s)$ is the unit normal vector to $\g$ at $A(s)=\g(s)$. Since the line $A_1A_2$ is not tangent to $\g_1$ at $A_1$, it
follows by item (i) of Proposition \ref{nnorm1}, that there is  a
$C^1$-smooth function $t_1(s)$, defined in a neighborhood $U$ of $0$, such that $A_1= \g_1(t_1(0))$ and
$\g_1(t_1(s))= \g(s)+ d_1(s)N(s)$ for every $s \in U$.
Here
$d_1(s):= \langle \g_1(t_1(s)) - \g(s), N(s)\rangle$ is
the {\it oriented distance} from $A(s)=\g(s)$ to $A_1(s)=\g_1(t_1(s))$ relative to the orientation
of the line $A_1A_2$ given by the vector $N(s)$. Set
$(T,N)=(T(0),N(0))$ and let $B_1$ be the intersection point of the
normal $n_{A_1}$ and the line $n_{C_A}$. This intersection point $B_1$ exists because $n_A$ is not tangent to $\g_1$ at $A_1$. We show next that
\begin{equation}\label{D1}
\overrightarrow{C_AB_1} =-\frac{d\,'_1(0)}{\kappa_{\g}(A)}T .
\end{equation}

To prove this, we differentiate the identity
$\g_1(t_1(s))= \g(s)+ d_1(s)N(s)$ and use the Frenet-Serret formula $N'(s)= -k(s)T(s)$ for plane
curves. This yields,  for the derivative at $s=0$, the expression

$$\frac{d}{ds}\g_1(t_1(s))|_{s=0}=
T+d\,'_1(0)N+d_1(0)N'(0)=T(1-d_1(0)\kappa_{\g}(A))+d\,'_1(0)N.$$

Set $\overrightarrow{C_AB_1} =\lambda T$. Then
$$ \overrightarrow{A_1B_1}= \overrightarrow{A_1C_A} +
\overrightarrow{C_AB_1}=
(\frac{1}{\kappa_{\g}(A)}-d_1(0))N+\lambda T$$ and, having in mind
that $C_A\neq A_1$ (i.e.
$\frac{1}{\kappa_{\g}(A)} \neq d_1(0)$), we conclude that $A_1B_1$ is perpendicular to
$\frac{d}{ds}\g_1(t_1(s))|_{s=0}$
if and only if $\lambda = -\frac{d'_1(0)}{\kappa(A)}$.

Let now $B_2$ be the intersection point of the normal $n_{A_2}$
and the line $n_{C_A}$. Then the same reasoning as above shows
that
\begin{equation}\label{D2}
\overrightarrow{C_AB_2} =-\frac{d\,'_2(0)}{\kappa_{\g}(A)}T ,
\end{equation}
where $d_2(s)$ is the oriented distance from $A(s)=\g(s)$ to
$A_2(s)=\g_2(\tau_2(s))$. The function
$\tau_2(s)$ here is obtained by item (i) of Proposition \ref{nnorm1} with the curve $\g_1$ replaced by the curve $\g_2$. It follows by (\ref{D1}) and (\ref{D2})
that
\begin{equation}\label{D3}
\overrightarrow{B_1B_2} =\frac{d\,'_1(0)-d\,'_2(0)}{\kappa_{\g}(A)}T.
\end{equation}
Let $d(s)$ be the oriented distance from $A_1(s)=\g_1(s)$ to
$A_2(s)=\g_2(s)$. Then $d(s)=d_2(s)-d_1(s)$ and
$d\,'_2(0)-d\,'_1(0)=d\,'(0)=0$ since $d(s)$ has a local extremum at
$s=0$. So, $B_1=B_2$ as desired.

In the sub-case (b) of case \emph{2} we can assume, without loss of generality, that $C_A= A_2$. This point belongs to
$n_{A_2}$ and to $n_{C_A}$. It remains to show that the line $A_1A_2$ is the normal to $\g_1$ at the point $A_1$. This will be derived from Lemma \ref{lem}. Indeed, as $A_1 \neq C_A$, it follows by item (ii) of Proposition \ref{nnorm1} and Corollary \ref{homo} that there exist an open interval
$V \subset \Delta_1$ containing $t_0$, and a  difeomorphism $s_1(t)$ defined in $V$ with values in
$\Delta$ such that $s_1(t_0) = 0$,
$A_1 =\g_1(t_0)$,
$A= \g(0)$, and for every
$t \in V$ the point $\g_1(t)$ belongs to the normal
 $n_{\g(s_1(t))}$. Note that the image $s_1(V)$ of the set $V$ is an open neighborhood of $0 $ in $\Delta$.

 Since the line $A_1A_2$ is not tangent to $\g_2$ at $A_2$, item (i) of Proposition \ref{nnorm1} implies that there exist an open interval $U \subset \Delta$ containing $0$, and a $C^1$-smooth function $\tau_2(s)$ defined in $ U$ with values in
$\Delta_2$ such that $\tau_2(0) = \tau_0$, $\g_2(\tau_0)=A_2$, and for every
$s \in U$ the point $\g_2(\tau_2(s))$ belongs to the normal $n_{\g(s)}$. Note that, by Proposition \ref{nnorm2} the derivative $\tau_2'(0) = 0$. We can assume, without restricting the generality, that
$s_1(V) \subset U$. This allows to consider the composed function $u(t) = \tau_2(s_1(t))$. Clearly, $u'(t_0) = 0$. Since $A_1A_2$ is a local extremum,
$d'(t)=0$ where $d(t)=|\g_1(t)\g_2(u(t))|$. Lemma \ref{lem} implies that the line $\g_1(t)\g_2(u(t)$ is the normal to $\g_1$ at the point $A_1$.


Let us now consider the case \emph{3} of item (i) from Theorem \ref{A}: the line $A_1A_2$ is tangent to $\g_1$ at $A_1$ and not tangent to
$\g_2$ at $A_2$.

%We will show that in this case the normals $n_{A_1}, n_{A_2}$ and the line $n_{C_A}$ intersect at one point.
If $A_1=C_A$, then
$n_{A_1}=n_{C_A}$ and  $n_{A_1}, n_{A_2}$ and $n_{C_A}$ intersect at one point because $A_1A_2$ is not tangent to $\g_2$ at $A_2$.

Let now $A_1\neq C_A$. Consider the function
$s_1(t)$ defined in item (ii) of Proposition \ref{nnorm1}. Note that (by Proposition
\ref{nnorm3}) $s_1'(t_0)=0$. Consider also the function $\tau_2(s)$ the existence of which is assured by item (i) of Proposition \ref{nnorm1}
(because $A_1A_2$ is not tangent to $\g_2$ at $A_2$.). Put $u(t):=\tau_2(s_1(t))$. Clearly, $u'(t_0)=0$. Lemma \ref{lem} implies that $A_1A_2$ is the normal to $\g_1$ at $A_1$. This is a contaradiction because $A_1A_2$ cannot be both a tangent and a normal to $\g_1$ at $A_1$ (because $\g_1'(t_0) \neq 0 $).

\medskip
It remains to prove the claim in item (ii) of  Theorem \ref{A} (when
$\kappa_{\g}(A)=0$).  There are, as above, three options for the line $A_1A_2$ which cover all possible cases:

\emph{1'. $A_1A_2$ is tangent to $\g_1$ at $A_1$ and  tangent to $\g_2$ at $A_2$.}

\emph{2'. $A_1A_2$ is not tangent to $\g_1$ at $A_1$ and not tangent to $\g_2$
at $A_2$. }

\emph{3'.}$ A_1A_2$ is tangent to $\g_1$ at $A_1$ and not tangent to $\g_2$
at $A_2$.

We proceed as in the proof of item (i).  The claim in the case \emph{1' } needs no proof. In \emph{3'} we reach a contradiction as in the case  \emph{3} of item (i). Therefore, it suffices to consider only the case \emph{2'} when the line
$A_1A_2$ is not tangent to $\g_1$ at $A_1$
and not tangent to $\g_2$ at $A_2$. We have seen that there exists an open set
$U$, $0 \in U \subset \Delta$, and two functions
$t_1(s)$ and $\tau_2(s)$ in $U$ with values in
$\Delta_1$ and $\Delta_2$ correspondingly,  such that $A_1=\g_1(t_1(0))$.
$A_2 = \g_2(\tau_2(0))$
and, for every $s \in U$, the points
$A_1(s):=\g_1(t_1(s)) $ and $A_2(s):=\g_2(\tau_2(s))$ belong to $n_{A(s)}$
where $A(s) = \g(s)$. Clearly,
$\g_1(t_1(s)) = \g(s) +d_1(s) N(s)$, where $d_1(s)$ is the oriented distance between
$A(s) = \g(s)$ and $A_1(s)=\g_1(t_1(s))$.  As
$N'(0)=-\kappa_{\g}(A)T=0$ we get for the derivative of $\g_1(t_1(s)) $ at $s=0$ the expression
$$\frac{d}{ds}\g_1(t_1(s))|_{s=0}=
T+d\,'_1(0)N.$$


Similarly, we have

$$\frac{d}{ds}\g_2(\tau_2(s))|_{s=0}=
T+d\,'_2(0)N$$
where $d_2(s)$ is the oriented distance between $A(s)$ and $A_2(s)$.
Since the triple $(A, A_1,A_2)$ is a local extremum
we have, as in item (i), that $d_1'(0)=d_2'(0)$. This implies that the tangent vector to
$\g_1(t)$ at $A_1$ and the tangent vector to $\g_2(\tau)$ at $A_2$ are colinear. It follows that the normals to the curves at $A_1$ and $A_2$ are parallel.
\end{proof}


\end{document}
